%% Manuscript for the Journal of Healthcare Informatics Research (Springer).
%% English version. Template: Springer Nature sn-jnl, reference style sn-basic
%% (numbered brackets, entries of the form Author, A.A.: Title. Journal vol,
%% pages (year)).
%%
%% Build: upload this single .tex into an Overleaf project that already contains
%% sn-jnl.cls and compile with pdfLaTeX. References are inline (thebibliography),
%% no .bib file needed.
%%
%% Keep manuscript_ru.md, manuscript_en.md and this file in sync when editing.

\documentclass[pdflatex,sn-basic,Numbered]{sn-jnl}

%%%% Standard template packages
\usepackage{graphicx}%
\usepackage{multirow}%
\usepackage{amsmath,amssymb,amsfonts}%
\usepackage{amsthm}%
\usepackage[title]{appendix}%
\usepackage{xcolor}%
\usepackage{textcomp}%
\usepackage{manyfoot}%
\usepackage{booktabs}%

\raggedbottom

\begin{document}

\title[Clinical Factors of Recurrent DKA]{Clinical Factors of Recurrent Diabetic Ketoacidosis: An Analysis of a Russian Cohort Using Statistics and Interpretable Machine Learning}

\author*[1]{\fnm{Oleg} \sur{Sadykhov}}\email{ovelsad@niuitmo.ru}

\author[1]{\fnm{Vasiliy} \sur{Leonenko}}\email{vnleonenko@itmo.ru}

\affil*[1]{\orgdiv{Artificial Intelligence Technologies Faculty}, \orgname{ITMO University}, \orgaddress{\street{49 Kronverksky Pr.}, \city{St. Petersburg}, \postcode{197101}, \country{Russia}}}

\abstract{\textbf{Purpose.} To identify clinical factors associated with a recurrent course of diabetic ketoacidosis (DKA) in a Russian cohort and to assess, using interpretable machine learning, how far these groups differ in clinical profile at admission.

\textbf{Methods.} A retrospective single-center study, predictors and outcome refer to the same hospitalization. Time to event was not recorded, so the model is diagnostic rather than prognostic. The outcome was established in 273 of 289 patients (recurrent 87, single 186). Groups were compared under a fixed statistical protocol with Benjamini-Hochberg correction. Independently, four interpretable models were fitted on 25 outcome-independent features without multicollinearity, the split preceding preprocessing. Feature contributions were estimated by SHAP and compared with the statistics. Two hypotheses were pre-specified, on alcohol intake and complex models outperforming logistic regression.

\textbf{Results.} Ten features were significantly associated with the outcome. Discrimination was moderate, ROC-AUC 0.71-0.78 by nested cross-validation and 0.72-0.80 on the held-out test. Five features entered the top-10 by SHAP in all four models, four statistically significant and the fifth (glucose at admission) found only multivariately. Under bootstrap, two of the five remained shared (89-92\% of replicates). The clinical hypothesis was confirmed, odds ratio 2.27 (95\% CI 1.23-4.20). The methodological hypothesis was not, DeLong found no significant differences between models.

\textbf{Conclusion.} The clinical profile separates the groups moderately, and the upper parts of the factor lists from two independent methodologies coincide.}

\keywords{diabetic ketoacidosis, recurrent course, risk factors, interpretable machine learning, SHAP, diagnostic model}

\maketitle

\section{Introduction}\label{sec1}

Diabetic ketoacidosis is a life-threatening complication of diabetes mellitus, characterized by hyperglycemia, metabolic acidosis and an elevated concentration of ketone bodies \cite{ref1}. Between 2003 and 2014 the United States recorded 1,760,101 primary DKA hospitalizations, and the number of discharges with a principal diagnosis rose from 118,808 in 2003 to 188,965 in 2014 \cite{ref2}. After the acute state is resolved, single cases may progress to repeated episodes, and a recurrent course carries substantially higher mortality \cite{ref3,ref4}.

Repeated DKA episodes are rarely explained by a single cause. Risk factors for repeated admissions include young age, psychiatric disorders, infections, non-adherence to therapy, alcohol and psychoactive substance abuse, and socioeconomic circumstances \cite{ref5,ref6,ref7}. The set of factors is heterogeneous, results across studies agree poorly, and for the Russian population the factors of a recurrent course have not been described. Despite decades of clinical experience, optimal DKA management strategies remain a subject of debate \cite{ref8}.

Closest to our task is a case-control observational study that matched patients with recurrent and single courses on age, diabetes duration, sex and ethnicity \cite{ref9}. After matching, few appreciable differences between the groups remained, and the main trigger of an episode in both groups was missed insulin injections. With 40 patient pairs such a result is hard to interpret unambiguously, it may mean either an absence of differences or a lack of power. The question of which patient characteristics are associated with a recurrent course remains open.

Machine learning allows the same task to be viewed from another angle. Published DKA models reach ROC-AUC of about 0.82-0.85 \cite{ref10,ref11}, but they address a different task, they predict the first episode or the fact of hospitalization within a set time window, they rely on large foreign electronic health record databases with a high proportion of missing values, and on other samples their metrics usually do not reproduce \cite{ref11}. It has also been shown \cite{ref12,ref13} that many models work as a black box, their conclusions are hard to interpret, which reduces clinicians' trust and becomes a barrier to adoption. Current work in diabetology applies prediction explanation methods, primarily SHAP and LIME \cite{ref14,ref15}, but for a recurrent course of DKA such tools have hardly been used.

In this work we identify the clinical factors associated with a recurrent course of DKA in a Russian cohort and assess how far a patient's clinical profile at admission can distinguish a recurrent course from a single episode. The key feature of the approach is that we obtain the set of factors along two independent paths, classical statistics under an agreed protocol and SHAP on top of interpretable machine learning models, after which we compare the results. Such a comparison shows to what extent the set of factors depends on the chosen analytical approach, which matters especially in a small sample.

We test two pre-specified hypotheses, the first of them clinical, concerning alcohol intake, which the literature describes as a factor of a recurrent course alongside psychoactive substance use, with excess alcohol named among the frequent triggers of an episode \cite{ref7}. The association is also supported by studies of repeated admissions \cite{ref9,ref16}. What remains is to find out whether it carries over to a Russian sample. The methodological hypothesis concerns model choice, a systematic review of clinical prediction models found no advantage of machine learning over logistic regression \cite{ref17}, yet this has not been tested on the task of recurrent DKA. The answer has a practical consequence, if complexity yields no gain, a simple and interpretable model is preferable.

\section{Methods}\label{sec2}

\subsection{Design and Data Source}\label{subsec2_1}

A retrospective single-center study on de-identified data with cross-sectional outcome ascertainment, both predictors and outcome refer to the same DKA hospitalization, with no observation period between them. The analysis included patients hospitalized with diabetic ketoacidosis between 2019 and 2025. Inclusion criteria were age 18 years or older and a confirmed diagnosis of DKA. The only exclusion criterion was an undetermined outcome. The initial sample comprised 289 patients, and the outcome was established in 273.

\subsection{Outcome}\label{subsec2_2}

The target variable was a recurrent course of DKA (0 for a single episode, 1 for a recurrent course). A patient was assigned to the recurrent group if more than one DKA episode had been recorded in total. The episode count combined the number of the patient's prior visits to this hospital with the same diagnosis from the medical records, the history at admission when the patient had not been observed here before, and the current hospitalization. A single course was defined as a first and only presentation.

The outcome was ascertained retrospectively, so at the time the predictors were measured it had in most cases already occurred. The data contain no fixed observation window, no dates of individual episodes and no time to event, and there is no censoring. Under the TRIPOD classification the model is diagnostic, it recognizes a state that has already occurred rather than predicting a future event. It does not admit direct comparison with models predicting hospitalization within 30, 90 or 180 days. The task addressed is the separation of a recurrent from a single course by the patient's clinical profile and the identification of factors associated with it. The consequences for interpretation are discussed in the limitations section.

\subsection{Hypotheses}\label{subsec2_3}

Both hypotheses were formulated in advance, before data analysis, on the basis of the literature.

\textbf{Clinical hypothesis.} H0: alcohol intake in the day before an episode is not associated with a recurrent course (odds ratio equals 1). H1: an association exists (odds ratio differs from 1). Test, the protocol test for a 2$\times$2 table, effect measure, odds ratio with a 95\% confidence interval, with the Benjamini-Hochberg correction for multiple comparisons.

\textbf{Methodological hypothesis.} H0: complex models (random forest, gradient boosting) and logistic regression do not differ in discrimination. H1: complex models outperform logistic regression by ROC-AUC. Test, the paired DeLong test for correlated ROC curves on the same patients, separately on the out-of-fold training part and on the held-out test set, with the Benjamini-Hochberg correction.

We put forward no other hypotheses. All other statistical comparisons are descriptive, and features with significant differences between groups are treated as an exploratory result requiring confirmation on an independent sample.

\subsection{Predictors}\label{subsec2_4}

We used indicators collected at hospitalization, demographic data (age, sex), history and therapy of diabetes (type, duration, age at onset, insulin delivery method, daily dose), laboratory indicators at admission (blood gases, electrolytes, creatinine, urea, glucose, HbA1c, lipid panel), complications (chronic kidney disease by stage and albuminuria, diabetic retinopathy and polyneuropathy) and a lifestyle factor (alcohol intake).

The models were trained on a single feature set, 25 indicators without marked multicollinearity. Age is almost linearly related to age at onset and diabetes duration (Pearson correlation about 0.99), so age at onset was removed as linearly dependent. The composition of the set was determined only by the mutual correlation of predictors and did not depend on the outcome, so feature selection introduces no bias into the quality estimate.

\subsection{Sample Size}\label{subsec2_5}

The sample size was not calculated in advance, the study is retrospective, so all available observations entered the analysis, 273 patients with an established outcome (recurrent course proportion 31.9\%, 87 events). With 87 events for 25 features there are few observations of the rare class, which raises the risk of overfitting, so we limited dimensionality and applied regularization and class weighting. Signs of overfitting were sought in three ways, by comparing estimates on nested cross-validation and the held-out test, by the out-of-bag estimate of the random forest, and by 95\% confidence intervals of the metrics from bootstrap.

\subsection{Missing Data}\label{subsec2_6}

The proportion of missing values per indicator ranged from a few percent to 58\% (HbA1c, sodium, potassium, lactate). We assumed the values were missing at random conditional on the observed variables, and accounted for their panel structure, where laboratory indicators go missing in groups. We compared imputation strategies, median and mode, k-nearest neighbors, and multiple imputation by chained equations. The imputer was trained only on the training part within folds, and on the test set missing values were filled from the nearest neighbors of the training part, without using the outcome or other test observations.

The working choice was k-nearest neighbors imputation with five neighbors. It outperformed the other strategies in model quality and did not appreciably distort the shape of the distributions, when comparing distributions before and after imputation by the Kolmogorov-Smirnov test no significant discrepancies were found. This check is indicative, since the compared samples are nested within one another, and it cannot be read as proof that the distributions coincide. The number of neighbors was compared separately (k from 3 to 11) by out-of-fold ROC-AUC on the training part, and k = 5 proved best.

\subsection{Statistical Analysis}\label{subsec2_7}

Normality was tested by the Shapiro-Wilk test for n below 50 and by the Kolmogorov-Smirnov test in the Lilliefors variant for n of 50 or more. Quantitative indicators with a normal distribution were described by mean and standard deviation with a 95\% confidence interval, and when departing from normality by median and quartiles. Groups were compared by the Student t-test, the Mann-Whitney U test or the Brunner-Munzel test depending on normality and equality of variances (Levene test). Proportions in 2$\times$2 tables were compared by the Pearson chi-square test when the minimum expected frequency exceeded 10 and by the Fisher exact test when the expected frequency was below 10. The confidence interval for proportions was computed by the Clopper-Pearson method, and the effect size by the odds ratio with a 95\% confidence interval. For multiple comparisons the Benjamini-Hochberg correction was applied. Areas under the ROC curves of two models on the same patients were compared by the DeLong test for correlated curves.

\subsection{Model Development}\label{subsec2_8}

The sample was split once, stratified, in an 80 to 20 ratio, training part 218 patients, test 55. The split was performed before any preprocessing, with a fixed random seed. All preprocessing (imputation, scaling, encoding of categorical features, balancing) was performed inside the pipeline and trained only on the training folds.

The choice of strategies and model families was made entirely on the training part, the held-out test was not consulted at this stage. We compared seven families (logistic regression, random forest, LightGBM, CatBoost, XGBoost, support vector machine, k-nearest neighbors) along the axes of imputation, balancing, categorical encoding and transformation of quantitative features, the comparison criterion being out-of-fold ROC-AUC on the training part.

Four families representing different model classes were taken forward, logistic regression as a linear model and the reference for the methodological hypothesis, random forest as a bagging method, and XGBoost and CatBoost as gradient boosting variants, the latter handling categorical features natively. The support vector machine and k-nearest neighbors lagged appreciably (out-of-fold ROC-AUC 0.677 and 0.631 against 0.713-0.780 for the others) and were not considered further. LightGBM was excluded as redundant with XGBoost, both implement gradient boosting over trees, and the difference between them (0.724 against 0.716) is small against the spread of estimates across folds.

Class imbalance was handled by class weighting, and synthetic methods (SMOTE, SMOTENC, BorderlineSMOTE, ADASYN) were compared separately. Hyperparameters were tuned by Bayesian search in nested cross-validation, the procedure quality was assessed by an outer five-fold cross-validation, while within each training fold the hyperparameters were tuned by a separate three-fold cross-validation.

\subsection{Performance Assessment}\label{subsec2_9}

The main metrics were chosen to be threshold-independent, ROC-AUC and PR-AUC with 95\% confidence intervals from bootstrap (2000 replicates). These answer the question of how far the clinical profile separates the groups and do not depend on an arbitrary threshold choice.

Probability calibration was checked by out-of-fold calibration curves and the Brier score, raw probabilities were compared with the Platt and isotonic regression methods, the application rule was set in advance, calibration is used only where it lowers the Brier score by at least 0.005.

In addition, for a clear view of the error structure, we report the confusion matrix. The cut-off point was chosen by a neutral criterion, the maximum of the Youden index on the out-of-fold probabilities of the training part, then fixed and applied to the held-out test, which was used once. This point does not serve as a clinical recommendation, since the outcome is ascertained retrospectively there is no screening task here, and the cut-off serves only to illustrate the ratio of type I errors (false positives) to type II errors (false negatives).

\subsection{Interpretation}\label{subsec2_10}

Feature contributions were estimated by SHAP for all four models, LinearSHAP for logistic regression, TreeSHAP for random forest and XGBoost, and the built-in SHAP value computation for CatBoost. The contributions of one-hot encoded categories were summed back to the source feature, and feature importance was measured by the mean absolute SHAP value.

Feature importance was examined at two levels, answering different questions.

The first level is agreement among the models. We checked whether the ranking of features coincides across models or depends on the choice of algorithm. We assessed this by the pairwise Spearman rank correlation between the full rankings, and by the intersection of the top lists, for which at thresholds of top-5, top-10 and top-15 we counted how many features enter the top by SHAP in all four models at once.

The dependence of this intersection on the sample was checked by bootstrap, 200 replicates, in each the training part was resampled with replacement, all four models were refitted with fixed hyperparameters, and SHAP was computed on the observations not drawn into the sample.

The second level is the comparison of machine learning with classical statistics. We checked whether the features shared across the top-10 of all models coincide with the list of indicators significant in the between-group comparison. The two lists were obtained independently, the results of the statistical comparison did not affect the set of features fed to the models. In addition, for each indicator from Table~\ref{tab1} we counted in how many of the four models it entered the top-10 by SHAP.

\subsection{Software and Reporting}\label{subsec2_11}

The analysis was performed in Python 3.11 with fixed library versions, and the code is open. Reporting follows the STROBE recommendations for observational studies, and for the part concerning model development the TRIPOD+AI recommendations were additionally taken into account.

\section{Results}\label{sec3}

\subsection{Participant Characteristics}\label{subsec3_1}

Of the initial cohort of 289 patients, 16 without an established outcome were excluded. The analytical sample comprised 273 patients, 87 with a recurrent course (31.9\%) and 186 with a single episode. The stratified split gave a training part of 218 patients (69 with a recurrent course, 149 with a single episode) and a test part of 55 (18 and 37 respectively).

Table~\ref{tab1} collects the indicators with a significant between-group difference after the Benjamini-Hochberg correction, 10 in total.

\begin{table}[h]
\caption{Features with a significant between-group difference (Benjamini-Hochberg correction, q below 0.05)}\label{tab1}%
\begin{tabular}{@{}lllll@{}}
\toprule
Indicator & Single (n=186) & Recurrent (n=87) & q & OR (95\% CI)\\
\midrule
Daily insulin dose, units & 34 [0; 45] & 44 [35; 54] & $<$0.0001 & -\\
No insulin therapy, \% & 28.9 & 1.1 & $<$0.0001 & -\\
HDL cholesterol, mmol/L & 1.09 [0.83; 1.38] & 1.26 [1.06; 1.58] & 0.0065 & -\\
Diabetes duration, years & 6 [0; 12.5] & 9 [4; 15] & 0.0139 & -\\
HbA1c, \% & 12.56 [10.28; 13.64] & 10.90 [8.91; 12.66] & 0.0196 & -\\
No retinopathy, \% & 65.4 & 48.1 & 0.0196 & -\\
No kidney damage (G0), \% & 56.3 & 25.8 & 0.0196 & -\\
No albuminuria (A0), \% & 52.7 & 27.8 & 0.0196 & -\\
Alcohol intake, \% & 15.3 & 29.1 & 0.0234 & 2.27 [1.23; 4.20]\\
Polyneuropathy, \% & 52.1 & 69.0 & 0.0281 & 2.05 [1.17; 3.57]\\
\botrule
\end{tabular}
\footnotetext{Quantitative indicators are given as median [Q1; Q3]. For categorical indicators the proportion of the stated level among known values in the group is given, missing values are not included in the denominator. The odds ratio is computed only for binary indicators.}
\end{table}

For the insulin delivery method, the main difference is driven not by the proportion on pump therapy but by the proportion of patients on no insulin therapy, 28.9\% in the single-episode group against 1.1\% in the recurrent group. This reflects a structural signal of disease onset, since a first episode by definition cannot be assigned to a recurrent course.

\subsection{Clinical Hypothesis: Alcohol Intake}\label{subsec3_2}

The feature value is known for 269 of 273 patients. Among patients with a recurrent course, 29.1\% (25 of 86) had consumed alcohol in the day before the episode, against 15.3\% (28 of 183) in the single-episode group. The minimum expected frequency in the 2$\times$2 table equals 16.9, so by protocol the Pearson chi-square test was applied, statistic 7.01 at one degree of freedom, p = 0.008. We reject the null hypothesis. The odds ratio was 2.27 (95\% CI 1.23-4.20), that is, alcohol intake the day before roughly doubles the odds of a recurrent course. After the Benjamini-Hochberg correction the association retains significance (q = 0.023). The direction of the effect agrees with the literature \cite{ref7,ref9,ref16}.

\subsection{Discrimination}\label{subsec3_3}

\begin{table}[h]
\caption{Discrimination by nested cross-validation}\label{tab2}%
\begin{tabular}{@{}lll@{}}
\toprule
Model & ROC-AUC & SD\\
\midrule
Random forest & 0.775 & 0.019\\
XGBoost & 0.748 & 0.038\\
CatBoost & 0.735 & 0.050\\
Logistic regression & 0.713 & 0.029\\
\botrule
\end{tabular}
\footnotetext{SD is the standard deviation across the outer folds of nested cross-validation.}
\end{table}

\begin{table}[h]
\caption{Discrimination on the held-out test}\label{tab3}%
\begin{tabular}{@{}ll@{}}
\toprule
Model & ROC-AUC (95\% CI)\\
\midrule
Logistic regression & 0.802 [0.672; 0.912]\\
XGBoost & 0.769 [0.630; 0.888]\\
CatBoost & 0.739 [0.587; 0.868]\\
Random forest & 0.724 [0.576; 0.851]\\
\botrule
\end{tabular}
\footnotetext{Confidence intervals were obtained by bootstrap (2000 replicates).}
\end{table}

Test estimates lie in the range 0.72-0.80, and nested cross-validation estimates in the range 0.71-0.78. The confidence intervals on the test are wide (0.26 in width on average) and overlap strongly, so the models cannot be ranked by the test result alone. Moreover the ranking reversed relative to nested cross-validation, the random forest, the leader on cross-validation, came last on the test, while logistic regression, last on cross-validation, gave the best result on the test. With 55 patients in the test part such a discrepancy is expected and itself serves as a warning against choosing a model by a single held-out sample.

PR-AUC on the training part holds in the range 0.52-0.55 against a random classifier baseline of about 0.32.

As a sign of overfitting we considered a systematic excess of the estimates obtained on the training part over the estimate on the held-out test beyond the 95\% confidence intervals. There is no such excess, for all four models the nested cross-validation estimate falls within the confidence interval of the test estimate. In addition, for the random forest two resampling schemes on the training part agree, the out-of-bag estimate 0.749 [0.683; 0.815] and the out-of-fold estimate 0.772 [0.710; 0.833] mutually cover each other with their confidence intervals. Neither check proves the absence of overfitting, but they show that the discrepancies between estimates do not exceed the statistical uncertainty at this sample size.

Pairwise comparison of the four models by the DeLong test with the Benjamini-Hochberg correction over 6 comparisons revealed no significant pair either on the out-of-fold training part or on the held-out test. The smallest values were obtained when comparing logistic regression with the complex models on the training part (p after correction 0.178). This does not prove the equality of the models, since at this sample size the power of the test is low, but neither does it give grounds to single out any model as the best.

\subsection{Methodological Hypothesis}\label{subsec3_4}

The pre-specified methodological hypothesis was tested separately (Table~\ref{tab4}). On the training part all three complex models are nominally above logistic regression (difference 0.04-0.05), but after correction the difference is not significant for any. On the held-out test the sign of the difference is reversed, logistic regression is above all three complex models by 0.03-0.08, and here too the differences are not significant.

\begin{table}[h]
\caption{Complex models versus logistic regression, DeLong test}\label{tab4}%
\begin{tabular}{@{}lllll@{}}
\toprule
Sample & Complex model & AUC complex & AUC logreg & p (BH)\\
\midrule
OOF & Random forest & 0.772 & 0.719 & 0.09\\
OOF & XGBoost & 0.761 & 0.719 & 0.09\\
OOF & CatBoost & 0.764 & 0.719 & 0.09\\
Test & Random forest & 0.724 & 0.802 & 0.50\\
Test & XGBoost & 0.769 & 0.802 & 0.55\\
Test & CatBoost & 0.739 & 0.802 & 0.50\\
\botrule
\end{tabular}
\footnotetext{OOF is out-of-fold predictions on the training part. The Benjamini-Hochberg correction is applied separately within each sample, over three comparisons.}
\end{table}

The null hypothesis could not be rejected, the superiority of complex models over logistic regression was not demonstrated on this sample. We stress the correct reading, this is not proof of the equivalence of the models, since with 218 patients on the training part and 55 on the test the power of the test is limited.

\subsection{Calibration and Error Structure}\label{subsec3_5}

Raw probabilities turned out to be better calibrated than any of the recomputed variants. The Brier score on out-of-fold predictions was 0.178-0.189 for raw probabilities, 0.185-0.192 after the Platt method, and 0.190-0.193 after isotonic regression. The rule set in advance (to apply calibration only if the Brier score drops by at least 0.005) was not met for any model, so raw probabilities are used throughout the further analysis.

We note a side effect, calibration also reduced discrimination (for the random forest from 0.772 to 0.744 after the Platt method). The sigmoid transformation is itself monotone and cannot change the order of probabilities, and with it the ROC-AUC. The drop is produced not by it but by the implementation of the procedure, the calibration wrapper refits the base model on cross-validation folds and averages their predictions, so the model itself changes, not only the probability scale. At this sample size this loss is noticeable.

The error structure on the held-out test at the Youden-index cut-off point is given in Table~\ref{tab5}.

\begin{table}[h]
\caption{Confusion matrix on the held-out test, Youden-index cut-off point}\label{tab5}%
\begin{tabular}{@{}llllllll@{}}
\toprule
Model & Threshold & TP & FP & FN & TN & Sens. & Spec.\\
\midrule
Logistic regression & 0.315 & 16 & 13 & 2 & 24 & 0.889 & 0.649\\
CatBoost & 0.250 & 14 & 15 & 4 & 22 & 0.778 & 0.595\\
XGBoost & 0.380 & 13 & 10 & 5 & 27 & 0.722 & 0.730\\
Random forest & 0.415 & 11 & 10 & 7 & 27 & 0.611 & 0.730\\
\botrule
\end{tabular}
\footnotetext{Test part, 18 patients with a recurrent course and 37 with a single episode. FP is a type I error, FN is a type II error.}
\end{table}

The matrix shows what is not visible from the ROC-AUC, at similar discrimination the models distribute errors between the two types differently. Logistic regression misses only 2 recurrent cases out of 18 but wrongly assigns 13 of 37 patients to that class. The random forest behaves the opposite way, fewer false alarms but 7 cases missed out of 18. XGBoost gives the most balanced picture. The spread underscores the instability of the cut-off point at this sample size, so we build substantive conclusions on threshold-independent metrics.

\subsection{Negative Results}\label{subsec3_6}

The effect of transforming quantitative features (manual clinical transformation and the Yeo-Johnson transformation against standardization only) was checked for all four models. There is no gain for any, for logistic regression the discrepancy between variants is 0.004 by ROC-AUC, for CatBoost 0.005, and for XGBoost the result coincides to the third decimal across all three variants (0.712). The last is expected, since decision trees are invariant to monotone transformations of features, and the logarithm together with the Yeo-Johnson transformation does not change the order of values, and hence the split points. The only noticeable discrepancy was given by the random forest (0.773 without transformation, 0.756 under manual, 0.775 under Yeo-Johnson), but even this has nothing to do with the transformation, the manual transformation variant is assembled from three feature blocks and changes the column order, while the random forest at a fixed seed selects random subsets of features by their indices.

Stacking with logistic regression as the meta-model over the out-of-fold predictions of the base models gave the best combination (random forest and XGBoost) with ROC-AUC 0.7699, which does not exceed the best single model (0.772). Adding base models did not improve quality.

Synthetic balancing also gave no gain. ROC-AUC for the synthetic methods was 0.715-0.725 (SMOTENC 0.716, BorderlineSMOTE 0.715, ADASYN 0.722, SMOTE over one-hot 0.725), whereas without balancing it was 0.729, and with class weighting 0.737. The Brier score meanwhile rises from 0.187 to 0.193-0.200, that is, calibration worsens. The best result was given by the non-synthetic class weighting, which was applied in all final models.

\subsection{Feature Importance}\label{subsec3_7}

The full feature rankings by mean absolute SHAP value agree moderately between models (Table~\ref{tab6}). As expected, agreement is highest between the random forest and XGBoost, both of which build tree ensembles. Agreement is lowest between CatBoost and the rest.

\begin{table}[h]
\caption{Pairwise Spearman rank correlation between SHAP rankings}\label{tab6}%
\begin{tabular}{@{}lllll@{}}
\toprule
 & CatBoost & Logreg & Forest & XGBoost\\
\midrule
CatBoost & 1.00 & 0.49 & 0.58 & 0.47\\
Logreg & 0.49 & 1.00 & 0.71 & 0.51\\
Forest & 0.58 & 0.71 & 1.00 & 0.91\\
XGBoost & 0.47 & 0.51 & 0.91 & 1.00\\
\botrule
\end{tabular}
\end{table}

The moderate agreement of the full rankings is explained by the ends of the lists, where the absolute SHAP values are small and the ranks unstable. The upper part of the lists coincides considerably more, and the size of the overlap depends on the cut-off threshold. In the top-5, two features are shared across all four models, insulin delivery method and daily insulin dose. In the top-10 there are five such features, HbA1c, HDL cholesterol and glucose at admission being added. In the top-15 there are eight. There is no distinguished number here, and the growth of the overlap as the threshold is relaxed is expected, there are only 25 features, and at a threshold covering the whole list all would be shared. What is informative, therefore, is not the growth but the composition of the overlap at a tight threshold, where a chance intersection is unlikely.

The overlap at the top-10 holds unevenly across the sample. Under bootstrap of the training part (200 replicates) insulin delivery method and daily insulin dose remain shared across all models in 92\% and 89\% of replicates, whereas HbA1c, HDL cholesterol and glucose at admission do so in 32-40\%. In individual models these three features enter the top-10 considerably more often, in 47-72\% of replicates. The gap between the two figures reflects a property of the criterion itself, the requirement of simultaneous presence in all four models is more sensitive to perturbation of the data than the ranking of each model separately. The same pair of features that is stable under resampling is already shared across all models at the tighter top-5 threshold. Without resampling the picture repeats, when SHAP is computed on the held-out test instead of the training part the rankings of individual models change little (Spearman 0.97-1.00), while the overlap at the top-10 shrinks from five features to four. Sample size also plays a part, at 218 observations the exclusion of a few patients noticeably changes the ranking.

The differences between models fall on features outside this overlap and reflect their construction. Logistic regression additionally raises electrolytes (sodium and potassium), the random forest and XGBoost raise diabetes duration and total cholesterol, and CatBoost raises the type of diabetes.

Comparison with classical statistics shows that of the five features shared across the top-10 of all models, four are also significant in the between-group comparison, insulin delivery method, daily insulin dose, HbA1c and HDL cholesterol. The fifth, glucose at admission, showed no significant between-group difference in the univariate comparison but entered the top-10 in all four models.

The full picture of correspondence is given in Table~\ref{tab7}. Indicators with the smallest q values enter the top-10 in all models, and as q grows the number of models declines. The agreement of the two approaches is strongest for the most significant indicators and gradually weakens.

\begin{table}[h]
\caption{Correspondence between significance in the between-group comparison and SHAP contribution}\label{tab7}%
\begin{tabular}{@{}lll@{}}
\toprule
Indicator & q & In top-10 SHAP, models of 4\\
\midrule
Insulin delivery method & $<$0.0001 & 4\\
Daily insulin dose & $<$0.0001 & 4\\
HDL cholesterol & 0.0065 & 4\\
HbA1c & 0.0196 & 4\\
Glucose at admission & not significant & 4\\
Kidney damage (CKD, G) & 0.0196 & 3\\
Urea at admission & not significant & 3\\
Total cholesterol & not significant & 3\\
Diabetes duration & 0.0139 & 2\\
Retinopathy & 0.0196 & 2\\
Albuminuria (CKD, A) & 0.0196 & 2\\
Polyneuropathy & 0.0281 & 1\\
Alcohol intake & 0.0234 & 0\\
\botrule
\end{tabular}
\footnotetext{Shown are indicators significant in the between-group comparison and indicators that entered the top-10 by SHAP in at least three models. The q value is taken from Table~\ref{tab1}.}
\end{table}

The discrepancies in the lower part of the table are expected, since the two quantities measure different things. The between-group test evaluates an isolated difference in a single indicator, whereas the mean absolute SHAP is the feature's contribution to the prediction, averaged over the whole cohort and computed in the presence of the remaining 24 variables. A binary feature with moderate prevalence receives a small value under such averaging, even if it is informative for its subgroup of patients.

\section{Discussion}\label{sec4}

We assessed how far a patient's clinical profile at admission can distinguish a recurrent course of DKA from a single episode, and identified the factors associated with the outcome. Discrimination holds in the range 0.71-0.78 on nested cross-validation and 0.72-0.80 on the held-out test.

One of the main results is the features shared by two independent approaches, the multivariate SHAP analysis and the univariate statistical comparison. Insulin delivery method, daily insulin dose, HbA1c, HDL cholesterol and glucose at admission enter the top-10 by SHAP in all four machine learning models, that is, they do not depend on the choice of algorithm. Four of these five features are also significant in the between-group comparison under the statistical protocol. This matters methodologically, classical statistics evaluates each feature in isolation, whereas SHAP evaluates its contribution within a multivariate model accounting for the remaining variables, so the approaches rest on different assumptions.

The approaches do not, however, coincide in everything, and these differences are of a regular nature. Glucose at admission entered the shared top-10 of the models, although it was not significant in the univariate between-group comparison, that is, the multivariate model singled out an indicator that on its own does not separate the groups. There are also the reverse cases, some significant indicators enter the top-10 in only some models (Table~\ref{tab7}). The statistical test and SHAP measure different things, so the discrepancy between them carries additional information rather than indicating an error. It is precisely this cross-comparison of two independent methodologies that forms the methodological side of the work, and on a small sample it helps to separate a signal common to the methods from the peculiarities of each of them.

The clinical hypothesis was confirmed, alcohol intake in the day before an episode is associated with a recurrent course, odds ratio 2.27 (95\% CI 1.23-4.20). This carries a risk factor known from the literature over to the Russian population \cite{ref7,ref9,ref16}. We note a limitation of interpretation, the design is cross-sectional, so this concerns an association rather than a proven causation.

The methodological hypothesis was not confirmed. The practical conclusion, at this sample size added model complexity does not pay off, and logistic regression is preferable as the simplest and most interpretable, with natural probability calibration. The random forest and XGBoost give the same quality through nonlinear interactions, and CatBoost is convenient for its native handling of categorical features, but none of them brings a gain. It is also telling that the model ranking on the held-out test came out reversed relative to the ranking on nested cross-validation, with 55 patients in the test part the differences between models are indistinguishable from random fluctuation. This negative result agrees with a systematic review \cite{ref17}, which on the material of other clinical tasks also found no advantage of machine learning over logistic regression.

A direct comparison of discrimination with published models is incorrect, since they address a different task. On a cohort of adults with type 1 diabetes a ROC-AUC of about 0.82 was obtained \cite{ref11}, up to 0.85 on a cohort of children \cite{ref10}, but both works predict the first episode or hospitalization within a set time window on large foreign databases.

\section{Limitations}\label{sec5}

The main limitation is the design. The outcome was ascertained retrospectively, so at the time the predictors were measured it had in most cases already occurred. The model is diagnostic, it does not predict a future event but separates already formed groups by clinical profile.

Diabetes duration differed between groups, a median of 9 years [4; 15] in the recurrent course against 6 years [0; 12.5] in the single episode (q = 0.0139). This difference may reflect not risk but the very definition of the outcome, the longer the disease lasts, the longer the period over which a repeat episode could be recorded, and the data contain no observation interval common to all patients. Building a prognostic model requires a prospective cohort with episode dates, observation time and censoring, which constitutes a natural continuation of the work.

The overlap of features across all models is only partially confirmed. Under bootstrap of the training part, insulin delivery method and daily insulin dose remain shared across all models in 92\% and 89\% of replicates, while HbA1c, HDL cholesterol and glucose at admission do so in 32-40\%, so the latter three should be regarded as requiring confirmation. The frequencies obtained are an upper bound, the hyperparameters were fixed under resampling, and the variability of their tuning is not built into the estimate. Resampling from a single cohort, moreover, says nothing about reproducibility on data from another center.

Separately we note a structural signal in the insulin delivery method feature. The proportion of patients on no insulin therapy was 28.9\% in the single-episode group against 1.1\% in the recurrent group, in our cohort the absence of insulin therapy at hospitalization almost always corresponds to disease onset, and a first episode by definition cannot be assigned to a recurrent course. The contribution of this feature therefore partly reproduces the definition of the outcome, and it cannot be read as an independent risk factor, although it is precisely this feature that, more often than the rest, turns out to be shared across all models.

The remaining limitations concern the data. The sample is single-center and small, with no external validation on an independent population. The proportion of missing values in part of the laboratory indicators is high, up to 58\%, and although imputation did not appreciably distort the distributions, uncertainty remains. A number of features named in the literature as leading for a recurrent course (psychiatric disorders, treatment adherence, socioeconomic circumstances) are absent from our data, and this limits the attainable discrimination.

\section{Conclusions}\label{sec6}

A patient's clinical profile at admission separates a recurrent from a single course of DKA with moderate ability, ROC-AUC 0.71-0.78 by nested cross-validation and 0.72-0.80 on the held-out test. Five features enter the top-10 by SHAP in all four models at once, insulin delivery method, daily insulin dose, HbA1c, HDL cholesterol and glucose at admission. Four of these five are also significant in the between-group comparison, that is, they are singled out along two independent methodological paths. Under bootstrap of the training part two of them remain shared across all models, insulin delivery method and daily insulin dose (92\% and 89\% of replicates), the other three in 32-40\%.

The clinical hypothesis was confirmed, alcohol intake in the day before an episode is associated with a recurrent course, odds ratio 2.27 (95\% CI 1.23-4.20). The methodological hypothesis was not confirmed, the superiority of complex models over logistic regression could not be demonstrated, and stacking and synthetic balancing likewise brought no advantage, which at this data volume argues in favor of simple interpretable models.

The model is diagnostic rather than prognostic over a time horizon. A prognostic model requires a prospective cohort with episode dates and observation time, together with external validation on data from another center.

\backmatter

\bmhead{Supplementary information}

The full result tables (cohort characteristics, comparison of preprocessing strategies, hyperparameter tuning, calibration, cut-off points and confusion matrices, pairwise DeLong comparisons, stacking, synthetic balancing, the frequencies with which features turn out to be shared across all models under bootstrap) and the SHAP figures are available in the open project repository.

\bmhead{Acknowledgements}

The authors thank the physicians who provided the de-identified clinical data.

\section*{Declarations}

\bmhead{Funding}

The study was carried out without external funding.

\bmhead{Conflict of interest}

The authors declare no conflict of interest relevant to the content of this article.

\bmhead{Ethics approval and consent to participate}

[To be completed after the ethics committee decision is obtained. Required form: the full committee name, protocol number and date, together with a statement of compliance with the Declaration of Helsinki of 1964 and its later amendments. For a retrospective study on de-identified data, a statement of exemption from the requirement of approval with the same details is admissible.]

\bmhead{Consent for publication}

Not applicable, the work was performed on de-identified data, individual patient information is not published.

\bmhead{Data availability}

Patient data cannot be published, as this would breach confidentiality. Access is possible on reasonable request to the corresponding author within ethical constraints.

\bmhead{Materials availability}

Not applicable.

\bmhead{Code availability}

The analysis code is open and available at \url{https://github.com/ovelsad/dka-recurrence-ml}, and includes fixed library versions and instructions for reproduction.

\bmhead{Author contributions}

O. Sadykhov: study conception, data preparation and analysis, statistical analysis, model development and validation, interpretation of results, preparation of the first draft of the manuscript. V. Leonenko: study conception and design, scientific supervision, analysis methodology, verification of results, critical revision and editing of the manuscript. Both authors read and approved the final version.

\begin{thebibliography}{17}

\bibitem{ref1} Benoit, S.R., Zhang, Y., Geiss, L.S., Gregg, E.W., Albright, A.: Trends in diabetic ketoacidosis hospitalizations and in-hospital mortality - United States, 2000-2014. MMWR Morb. Mortal. Wkly. Rep. \textbf{67}(12), 362-365 (2018). \url{https://doi.org/10.15585/mmwr.mm6712a3}

\bibitem{ref2} Desai, D., Mehta, D., Mathias, P., Menon, G., Schubart, U.K.: Health care utilization and burden of diabetic ketoacidosis in the U.S. over the past decade: a nationwide analysis. Diabetes Care \textbf{41}(8), 1631-1638 (2018). \url{https://doi.org/10.2337/dc17-1379}

\bibitem{ref3} Santos, S.S., Ramaldes, L.A.L., Dualib, P.M., Gabbay, M.A.L., Sa, J.R., Dib, S.A.: Increased risk of death following recurrent ketoacidosis admissions: a Brazilian cohort study of young adults with type 1 diabetes. Diabetol. Metab. Syndr. \textbf{15}(1), 85 (2023). \url{https://doi.org/10.1186/s13098-023-01054-5}

\bibitem{ref4} Gibb, F.W., Teoh, W.L., Graham, J., Lockman, K.A.: Risk of death following admission to a UK hospital with diabetic ketoacidosis. Diabetologia \textbf{59}(10), 2082-2087 (2016). \url{https://doi.org/10.1007/s00125-016-4034-0}

\bibitem{ref5} Mohler, R., Lotharius, K., Moothedan, E., Goguen, J., Bandi, R., Beaton, R., Knecht, M., Mejia, M.C., Khoury, M., Sacca, L.: Factors contributing to diabetic ketoacidosis readmission in hospital settings in the United States: a scoping review. J. Diabetes Complications \textbf{38}(10), 108835 (2024). \url{https://doi.org/10.1016/j.jdiacomp.2024.108835}

\bibitem{ref6} Dhatariya, K.K., Glaser, N.S., Codner, E., Umpierrez, G.E.: Diabetic ketoacidosis. Nat. Rev. Dis. Primers \textbf{6}(1), 40 (2020). \url{https://doi.org/10.1038/s41572-020-0165-1}

\bibitem{ref7} Brandstaetter, E., Bartal, C., Sagy, I., Jotkowitz, A., Barski, L.: Recurrent diabetic ketoacidosis. Arch. Endocrinol. Metab. \textbf{63}(5), 531-535 (2019). \url{https://doi.org/10.20945/2359-3997000000158}

\bibitem{ref8} Kanzwl, S.M.O.S., Alhajri, A.H.M., Mohammed, Y.J.A., et al.: A comparative effectiveness of intravenous fluids and insulin regimens in the acute management of diabetic ketoacidosis (DKA) and hypoglycemia: a systematic review. Cureus \textbf{17}(10), e94902 (2025). \url{https://doi.org/10.7759/cureus.94902}

\bibitem{ref9} Cooper, H., Tekiteki, A., Khanolkar, M., Braatvedt, G.: Risk factors for recurrent admissions with diabetic ketoacidosis: a case-control observational study. Diabet. Med. \textbf{33}(4), 523-528 (2016). \url{https://doi.org/10.1111/dme.13004}

\bibitem{ref10} Williams, D.D., Ferro, D., Mullaney, C., Skrabonja, L., Barnes, M.S., Patton, S.R., Lockee, B., Tallon, E.M., Vandervelden, C.A., Schweisberger, C., Mehta, S., McDonough, R., Lind, M., D'Avolio, L., Clements, M.A.: An ``all-data-on-hand'' deep learning model to predict hospitalization for diabetic ketoacidosis in youth with type 1 diabetes: development and validation study. JMIR Diabetes \textbf{8}, e47592 (2023). \url{https://doi.org/10.2196/47592}

\bibitem{ref11} Li, L., Lee, C.C., Zhou, F.L., Molony, C., Doder, Z., Zalmover, E., Sharma, K., Juhaeri, J., Wu, C.: Performance assessment of different machine learning approaches in predicting diabetic ketoacidosis in adults with type 1 diabetes using electronic health records data. Pharmacoepidemiol. Drug Saf. \textbf{30}(5), 610-618 (2021). \url{https://doi.org/10.1002/pds.5199}

\bibitem{ref12} Rudin, C.: Stop explaining black box machine learning models for high stakes decisions and use interpretable models instead. Nat. Mach. Intell. \textbf{1}(5), 206-215 (2019). \url{https://doi.org/10.1038/s42256-019-0048-x}

\bibitem{ref13} Stiglic, G., Kocbek, P., Fijacko, N., Zitnik, M., Verbert, K., Cilar, L.: Interpretability of machine learning-based prediction models in healthcare. WIREs Data Min. Knowl. Discov. \textbf{10}(5), e1379 (2020). \url{https://doi.org/10.1002/widm.1379}

\bibitem{ref14} Ahmed, S., Kaiser, M.S., Hossain, M.S., Andersson, K.: A comparative analysis of LIME and SHAP interpreters with explainable ML-based diabetes predictions. IEEE Access \textbf{13}, 37370-37388 (2025). \url{https://doi.org/10.1109/ACCESS.2024.3422319}

\bibitem{ref15} Emi-Johnson, O.G., Nkrumah, K.J.: Predicting 30-day hospital readmission in patients with diabetes using machine learning on electronic health record data. Cureus \textbf{17}(4), e82437 (2025). \url{https://doi.org/10.7759/cureus.82437}

\bibitem{ref16} Bradford, A.L., Crider, C.C., Xu, X., Naqvi, S.H.: Predictors of recurrent hospital admission for patients presenting with diabetic ketoacidosis and hyperglycemic hyperosmolar state. J. Clin. Med. Res. \textbf{9}(1), 35-39 (2017). \url{https://doi.org/10.14740/jocmr2792w}

\bibitem{ref17} Christodoulou, E., Ma, J., Collins, G.S., Steyerberg, E.W., Verbakel, J.Y., Van Calster, B.: A systematic review shows no performance benefit of machine learning over logistic regression for clinical prediction models. J. Clin. Epidemiol. \textbf{110}, 12-22 (2019). \url{https://doi.org/10.1016/j.jclinepi.2019.02.004}

\end{thebibliography}

\end{document}
